\section{Conclusion}

This draft summarizes a baseline-first limitations audit for finite
superelliptic quantum dots on a square tight-binding lattice. The system remains
useful as a controlled benchmark: the Hamiltonian convention is clear, the
geometry family is compact, and the direct Kwant calculations provide a
consistent reference for the tested finite systems. The audit shows, however,
that several attractive objectives fail to support positive claims once simple
baselines and methodological checks are applied.

For the closed spectra, the physics-informed Ridge model remains the preferred
conservative surrogate, while a small MLP does not meet the robust structured
validation criterion. Q and S inverse-screening objectives do not beat their
isotropy, symmetry, or anisotropy baselines. The magnetic ranking-crossover
check is explained by symmetry-lifting baselines. The embedded-mask
finite-difference reference is not adequate as a high-precision continuum truth
for shape-sensitive residuals, and the TB-only self-scaling fallback remains
inconclusive. In the open-transport preflight, contact width dominates the
tested shape comparison.

These results do not prove a universal no-go theorem for geometry-based quantum
design, machine learning, continuum references, or transport objectives. They
show that this particular superellipse metric search does not justify a
positive design rule. Future work should therefore change the physical question
or the experimental/numerical protocol rather than continue minor variants of
the same objective. A stronger future direction would need to define the
baseline ladder in advance, control reference uncertainty and contact geometry,
and make positive claims only after the surviving signal is direct, stable, and
not absorbed by simpler physics.
